\section{A defect found during revision, and its repair}
\label{sec:supp_defect}

We record this because the released code contains a switch that restores the earlier behavior, and a
reader who finds that switch deserves to know what it is for.

\subsection{What was wrong}
In the version of this work that was first submitted, the allocation was formed by adding the rule
contributions to a base level and then clipping the result to the unit interval. The positive
contributions could reach $0.884$ and $0.762$ on top of a base of $0.818$, so the value arriving at
the clip could be as large as $2.46$. The clip was therefore active at every pixel of every image and
the allocation was the constant one.

Two consequences follow. The rule layer had no influence on any output, since multiplying by one
changes nothing. And the derivative of a clip at its bound is zero, so no gradient reached the rule
parameters and they stayed at their initial values for the whole of training. We confirmed both
directly, by reading the allocation map out of the trained model and by comparing every rule
parameter in the saved weights against its initializer after 2082 optimizer steps.

\subsection{The repair}
The allocation is now a bounded weighted vote. Writing $r_k \in [0,1]$ for the truth value of rule
$k$ and $w_k \in (0,1)$ for its learned weight,
\begin{equation}
\begin{split}
a = {}& B_{\min} + (B_{\max}-B_{\min})\,\sigma(\beta) \\
     & - C_1 w_1 r_1 + C_2 w_2 r_2 + C_3 w_3 r_3 \\
     & + C_4 w_4 r_4 - C_5 w_5 r_5 ,
\end{split}
\end{equation}
with $B_{\min} = 0.25$, $B_{\max} = 0.50$ and budgets $C_1 = 0.15$, $C_2 = 0.25$, $C_3 = 0.15$,
$C_4 = 0.10$ and $C_5 = 0.10$. The budgets satisfy
\begin{equation}
B_{\max} + C_2 + C_3 + C_4 = 1, \qquad B_{\min} - C_1 - C_5 = 0 ,
\end{equation}
so the value lies in the unit interval before any clipping, for every input and every parameter
setting. Clipping is therefore never active, saturation is impossible rather than unlikely, and the
gradient reaches every rule parameter.

The repair also tightens the theory. Because the allocation is affine in the rule truth values with
signed coefficients and no clipping intervenes, the Lipschitz constant of Theorem~1 of the main paper
is attained rather than merely an upper bound.

This is also what the trained model shows and not only a property of the formula. Measured on the
released weights, the allocation now spans $\AllocMin$ to $\AllocMax$ with mean $\AllocMean$ and
standard deviation $\AllocStd$, and $\AllocClamp$ percent of pixels sit at either end of the clamp,
against the whole image before the repair, direct confirmation that no pixel saturates in practice.

\subsection{Reproducing the earlier behavior}
Setting the environment variable \texttt{LEGACY\_SATURATING\_ALLOCATION} to $1$ restores the original
computation exactly. We verified this against a copy of the pre revision source on twenty random
inputs, comparing the allocation maps for equality and the rule traces for agreement. Both matched.
The earlier numbers can therefore be recovered from the released code, which is why the switch exists.
